% sec-11-references-backmatter.tex - References and Back Matter (final)
% Unnumbered back matter so the numbered IND sections end at 11 (Relevant
% Information), matching the ReGARDD IND Table-of-Contents numbering.
\section*{References and Back Matter}
\addcontentsline{toc}{section}{References and Back Matter}

This back matter closes the Initial Investigational New Drug Application v1.0 for
the LLM-Directed PDAC Robotic Daraxonrasib Phase 1 study, assembled within
repository v4.3.0. The complete reference list renders below through the
\texttt{ieeetr} bibliography style from \texttt{references.bib}. Every entry
carries a clickable URL and, where a digital object identifier is assigned, a
clickable DOI URL with the DOI text shown in full; long links break on any
character so that no link runs off the right margin of the text block. The
reference base supporting this application is intentionally comprehensive and
spans seven evidentiary strands: the FDA Phase 1 and combination drug-device
regulatory guidances \cite{phase1guidance}; the daraxonrasib (RMC-6236) clinical
and mechanistic literature underpinning the RASolute 302 and RASolve 301 programs
and the June 2025 Breakthrough Therapy designation \cite{daraxwhipple}; the
21 CFR Part 312 (IND) and Part 50 (informed consent) and ICH Physical Artificial
Intelligence regulatory adaptations \cite{cfr312paiuni,cfr050paiuni,ichpaistduni};
the five author computational works dated 2025 to 2026 (the 60-second PDAC Whipple
plus daraxonrasib simulation, the artificial intelligence digital-twin PDAC
simulation, the quantitative systems pharmacology metastatic-PDAC simulation with
verification, validation, and uncertainty quantification, the 100,000-patient in
silico Phase 3 five-arm PDAC triplicate, and the end-to-end PDAC digital-twin trial
proposals) \cite{chatgpt100kp,qspmetpancre,fdadigtwinpc,pdacdigtwinp}; the H.R. 9510
Verification Before Generation in Physical AI Oncology Trials Act of 2026 and the
national platform and clinical-trust frameworks \cite{congress9510,clintrustpai,natlpaiplatf};
the sixteen-large-language-model code-generation evaluation against the FDA Adverse
Event Reporting System \cite{llmcomp16fda}; and the proposed clinical research
protocol v1.0.0 (DOI 10.5281/zenodo.20780121) that this IND operationalizes.

The back matter elements that follow document authorship and contributor
attribution, the open-weights build model used to assemble this application, the
ethical and independence disclosures that bound the entire repository, the data
availability and licensing terms, and the recommended citation. Table -.1 records
the provenance of each major content strand so that a reviewer can trace every
quantitative claim in IND Sections 3 through 11 back to its grounding source.
Table -.2 then maps the back matter disclosure elements to the controlling federal
authorities under 21 CFR, and Table -.3 enumerates the four build-stage source
trees that constitute the archived, reproducible record of this application.

\needspace{12\baselineskip}
\begin{table}[htbp]
\centering
\begin{tabularx}{\textwidth}{>{\raggedright\arraybackslash}p{3.6cm} Y >{\raggedright\arraybackslash}p{2.3cm}}
\toprule
\textbf{Evidentiary strand} & \textbf{Grounding source and key facts} & \textbf{Supports IND sections} \\
\midrule
Daraxonrasib mechanism and class & Oral RAS(ON) multi-selective (pan-RAS) inhibitor binding GTP-bound mutant RAS in complex with cyclophilin A; once-daily dosing \cite{daraxwhipple} & 3.1, 8.1 \\
Dose-escalation design & 3+3 escalation across DL1 160 mg, DL2 220 mg, DL3 300 mg once daily; 28-day DLT window; RP2D at or below MTD \cite{phase1guidance} & 4, 6.1 \\
Robotic Whipple device & On-premises LLM-directed eight-arm system, 56 degrees of freedom, 640 sensor channels, force caps 3 N per arm and 18 N cumulative \cite{onpremwhippl} & 5, 7 \\
Verification before generation & Ten-gate ACCEPT/BLOCK/ESCALATE suite; LLM as second-opinion advisory oracle only; full autonomy prohibited under 21 CFR \S 312.21(e) \cite{congress9510,clintrustpai} & 5, 10, 11 \\
Perioperative pause-and-restart & ISGPS fistula grade plus serum trough versus 0.5 ng/mL; 32-iteration sweep yielding T+7d 29/32, T+14d 3/32, T+21d 0/32 \cite{paisitedocpk} & 6.1, 9 \\
Computational evidence base & 100,000-patient in silico Phase 3, QSP metastatic-PDAC simulation, and digital-twin trials \cite{chatgpt100kp,qspmetpancre,fdadigtwinpc,pdacdigtwinp} & 8, 9 \\
Regulatory framework & 21 CFR Part 312 and Part 50 Physical AI adaptations; ICH harmonization \cite{cfr312paiuni,cfr050paiuni,ichpaistduni} & 1, 6, 10, 11 \\
\bottomrule
\end{tabularx}
\vspace*{-0.4cm}
\figcaption{Table -.1. Provenance map linking evidentiary strand to its
grounding reference and to the IND sections it supports.}
\end{table}

\bmhead{Acknowledgments}
This Investigational New Drug Application was assembled with Claude Code (Opus 4.8
Ultracode build, 1M context) as the build model, operating under the single-prompt
mermaid to draft to full to final pipeline described in IND Section 11 and committing
each stage in real time. At each stage the build model expanded a grayscale Mermaid
specification into a draft section scaffold, then into a full regulatory draft, and
finally into this senior-author polished form, with automated commits and pull
requests opened for human review after every stage so that the provenance of every
sentence and every quantitative value remained auditable. 

The author thanks
ChatGPT 5.5 (Thinking Extended) and Gemini 3.1 Pro for research-source assistance
and for manual editing support; these systems were used only for source retrieval
and editorial review and exercised no control over the clinical design, the safety
architecture, or the regulatory determinations recorded herein. This v1.0
application was edited and approved by Kevin Kawchak and added to repository
v4.3.0 under
\href{https://github.com/kevinkawchak/physical-ai-oncology-trials/tree/main/trial-ind}{kevinkawchak/physical-ai-oncology-trials/tree/main/trial-ind}.
The repository contributors are @kevinkawchak, @claude, @google-gemini, and
@openai. No funding body, sponsor, or regulatory authority directed, reviewed, or
endorsed the preparation of this document.

\bmhead{Author and ORCID}
Kevin Kawchak, Chief Executive Officer, ChemicalQDevice, San Diego, California.
ORCID iD:
\href{https://orcid.org/0009-0007-5457-8667}{https://orcid.org/0009-0007-5457-8667}.
Correspondence: \href{mailto:kevink@chemicalqdevice.com}{kevink@chemicalqdevice.com}.
In the proposed study the author would serve as Sponsor-Investigator on behalf of
ChemicalQDevice under the combined obligations of a sponsor and an investigator set
out in 21 CFR \S 312.3(b), assuming the sponsor duties of 21 CFR \S\S 312.50 through
312.59 and the investigator duties of 21 CFR \S\S 312.60 through 312.69
simultaneously. The resulting financial interest is disclosed in full under
21 CFR Part 54 and on FDA Form 3455, and is mitigated by the independent three-tier
oversight structure detailed in the proposed clinical research and additional
information sections: a Data Safety Monitoring Board with binding halt authority at
each cohort boundary and on any triggered review, an Independent Safety Monitor with
real-time per-procedure stop authority, and a Physical AI Safety Review Committee
convening on a cadence not exceeding 90 days. Because the Sponsor-Investigator does
not sit on any of these bodies and cannot override a stop or halt determination, the
conflict of interest is structurally separated from the decisions that protect
enrolled subjects, consistent with the financial-interest mitigation expectations
under 21 CFR Part 54.

\bmhead{Ethical Disclosures}
This is an independent research and regulatory-authoring demonstration. No real
patients were enrolled, no protected health information was used, and no
Institutional Review Board approval was sought or obtained. The exemplar case
PAT-PDAC-0001 (a synthetic 68-year-old with a 3.4 cm head-of-pancreas lesion
abutting the superior mesenteric vein, KRAS G12D, microsatellite stable, CA 19-9
of 412 U/mL, ECOG 1, after four cycles of neoadjuvant mFOLFIRINOX) and all
quantitative simulation outputs, including the 32-iteration pause-and-restart sweep
and the 100-tick vascular no-fly verdict distribution, are synthetic and
illustrative and were generated in silico for the sole purpose of exercising the
regulatory and safety logic of this application. This document is not an active
Investigational New Drug Application or Investigational Device Exemption, is not a
submitted or approved protocol, and does not constitute medical, surgical, or
regulatory advice. 

It is not endorsed by the United States Food and Drug
Administration, the National Institutes of Health, the Department of Health and
Human Services, any Institutional Review Board, the International Council for
Harmonisation, or any sponsor. Any future first-in-human conduct of the design
described here would require a separately submitted IND under 21 CFR Part 312 and a
separately submitted IDE under 21 CFR Part 812, a 30-day FDA review period that
elapses without the imposition of a clinical hold under 21 CFR \S 312.42, and
prospective IRB approval under 21 CFR Part 56 together with informed consent
obtained under 21 CFR Part 50 and ICH E6(R3). The staged-autonomy framework that
constrains this Stage 1 study to clinician-controlled teleoperation and supervised
operation, with the LLM acting only as a second-opinion advisory oracle and full
autonomy prohibited under 21 CFR \S 312.21(e), reflects a deliberate ethical choice
to keep a qualified human surgeon in unambiguous control of every committed action.

\needspace{12\baselineskip}
\begin{table}[htbp]
\centering
\begin{tabularx}{\textwidth}{>{\raggedright\arraybackslash}p{3.5cm} Y >{\raggedright\arraybackslash}p{3.4cm}}
\toprule
\textbf{Back matter element} & \textbf{Disclosure content} & \textbf{Controlling authority} \\
\midrule
Sponsor-Investigator role & Combined sponsor and investigator obligations assumed by a single party & 21 CFR \S 312.3(b); \S\S 312.50 to 312.69 \\
Financial interest & Sponsor-Investigator equity interest disclosed and mitigated by independent oversight & 21 CFR Part 54; FDA Form 3455 \\
Human subjects status & No enrollment, no PHI, no IRB approval; synthetic exemplar only & 21 CFR Part 50; Part 56 \\
Future conduct gate & Separate IND plus IDE, 30-day review with no clinical hold, prospective IRB approval & 21 CFR \S\S 312.40 to 312.42; Part 812 \\
Autonomy limit & Stage 1 teleoperation and supervised operation; advisory oracle only & 21 CFR \S 312.21(e) \\
Data and records & CC BY 4.0 archive; reproducible build sources; no proprietary data & 21 CFR Part 11 (electronic records) \\
\bottomrule
\end{tabularx}
\vspace*{-0.4cm}
\figcaption{Table -.2. Back matter disclosure elements mapped to the controlling
federal authorities that govern each element.}
\end{table}

\bmhead{Data Availability}
All application sources, spanning the \texttt{mermaid}, \texttt{draft-ind},
\texttt{full-ind}, and \texttt{final-ind} build stages, are available in the public
GitHub repository and archived on Zenodo under the Creative Commons Attribution 4.0
International (CC BY 4.0) license at DOI
\href{https://doi.org/10.5281/zenodo.21097442}{10.5281/zenodo.21097442}. The four
build stages constitute a complete and reproducible authoring record: the mermaid
stage holds the grayscale figure specifications, the draft-ind stage holds the
section scaffolds, the full-ind stage holds the regulatory drafts, and the final-ind
stage holds this senior-author polished application. 

The underlying proposed clinical
research protocol v1.0.0 is archived separately at DOI
\href{https://doi.org/10.5281/zenodo.20780121}{10.5281/zenodo.20780121}, and the five
supporting computational works are archived at their own DOIs as cited in the
bibliography below. No proprietary, embargoed, or patient-identifiable data are
contained in or required to reproduce this application; every numeric value can be
regenerated from the deterministic, hash-pinned simulation seeds recorded in the
source tree, consistent with the electronic-records integrity expectations of
21 CFR Part 11.

\needspace{11\baselineskip}
\begin{table}[htbp]
\centering
\begin{tabularx}{\textwidth}{>{\raggedright\arraybackslash}p{2.8cm} Y >{\raggedright\arraybackslash}p{4.2cm}}
\toprule
\textbf{Build stage} & \textbf{Contents and role in the reproducible record} & \textbf{Archive} \\
\midrule
\texttt{mermaid} & Grayscale Mermaid figure specifications (8-tone classDef ramp) that map 1:1 to the TikZ figures & Zenodo CC BY 4.0; DOI 10.5281/zenodo.21097442 \\
\texttt{draft-ind} & Section scaffolds with drafting instructions and the ReGARDD section file map & Same repository and DOI \\
\texttt{full-ind} & Full regulatory drafts with complete tables, figures, and citations & Same repository and DOI \\
\texttt{final-ind} & Senior-author polished application; full-width ragged-right tables; verified figures & Same repository and DOI \\
\bottomrule
\end{tabularx}
\vspace*{-0.4cm}
\figcaption{Table -.3. Four build-stage source trees that form
archived, reproducible record of IND v1.0 under repository v4.3.0.}
\end{table}

\bmhead{Rights and Permissions}
This article is distributed under the terms of the Creative Commons Attribution 4.0
International License (CC BY 4.0), which permits unrestricted use, distribution, and
reproduction in any medium, provided the original author and source are properly
credited, a link to the Creative Commons license is provided, and any modifications
made are indicated. The license imposes no additional restrictions and does not
limit fair use, the rights of the author, or other applicable rights. To view a
copy of this license, visit
\url{https://creativecommons.org/licenses/by/4.0/}.

\bmhead{Cite This Article}
Kawchak K. Investigational New Drug Application - Daraxonrasib, Phase 1, LLM-Directed Robotic Whipple in KRAS-Mutated PDAC. Zenodo. 2026;
\href{https://doi.org/10.5281/zenodo.21097442}{10.5281/zenodo.21097442}. IND v1.0,
repository v4.3.0.

% \clearpage

\vspace{1.1cm}

% Bibliography (ieeetr); \nocite{*} emits the full assembled reference list.
\phantomsection\addcontentsline{toc}{section}{References}
\begingroup\raggedright\nocite{*}\bibliographystyle{ieeetr}\bibliography{references}\endgroup
